
\begin{tcolorbox}[promptbox={\textsc{ClaimExtract}: Extract verifiable claims from caption}]
Extract ONLY verifiable claims explicitly stated in the caption, limited to: color, position, and quantity.

Guidelines:
- Only extract claims that are explicit and directly stated in the caption.
- Do NOT infer unstated attributes or any information not present in the caption.
- Each claim MUST be grounded in an exact text span from the caption as "evidence".
- Each claim MUST belong to exactly one of three types: color, position, or quantity.
- Each claim MUST clearly associate the attribute with a specific object.
- If no valid claim exists, output [].

[Color] Include explicit colors (e.g., "red", "light gray"). Example: evidence: "a black cat" → claim: "The cat is black." \newline
[Position] Allowed: left, right, center, top, bottom only. Format: "The [object] is at the [position] of the image." \newline
[Quantity] Explicit numbers only; exclude vague quantifiers (e.g., "several", "many").

\medskip
Input: [Caption]: \{response\} \newline
Output: [\{"evidence": "...", "claim": "...", "type": "color/position/quantity"\}, ...]
\end{tcolorbox}

\begin{tcolorbox}[promptbox={\textsc{ClaimToQA}: Convert claims into question–answer pairs}]
Convert claims into Question-Answer pairs using templates.

Guidelines:
- Do NOT rewrite or reinterpret the claim.
- The question MUST be directly derived from the claim.
- The answer MUST be a substring or direct normalization of the claim.

[Color] Q: What color is the [OBJECT]? \ A: [color phrase] \newline
[Position] Q: Where is the [OBJECT] in the image? \ A: left / right / top / bottom / center \newline
[Quantity] Q: How many [OBJECT] are in the image? \ A: [number]

\medskip
Input: [Claims]: \{caption\_claim\} \newline
Output: [\{"claim": "...", "question": "...", "answer": "..."\}, ...]
\end{tcolorbox}

\begin{tcolorbox}[promptbox={\textsc{ImageQA}: Answer questions grounded in the image}]
Answer the questions based on the image.

Guidelines:
- Answer each question based only on the image.
- Color questions: answer with a color word or phrase only.
- Position questions: answer with ONLY one of: left, right, center, top, bottom.
- Quantity questions: answer with a number word or numeral only.
- If the answer cannot be determined reliably from the image, answer "unknown".

\medskip
Input: [Questions]: \{question\} \newline
Output: [\{"question": "...", "answer": "..."\}, ...]
\end{tcolorbox}

\begin{tcolorbox}[promptbox={\textsc{AnswerCheck}: Compare caption answer vs.\ image answer}]
Compare Answer1 (from caption) and Answer2 (from image).

Rules:
- Output 1 if they mean the same thing; output 0 otherwise.
- Output 0 if Answer2 is "unknown" or either answer is empty.

Matching guidelines:
- Ignore format differences: "3" == "three"
- Allow simple variants: "blue" == "dark blue"
- Allow position variants: "left" == "on the left"
- Otherwise, be strict.

\medskip
Input: [Answer Pairs]: \{check\_list\} \newline
Output: [\{"answer1": "...", "answer2": "...", "is\_correct": 1/0\}, ...]
\end{tcolorbox}
